\documentclass[12pt]{article}
\usepackage[utf8]{inputenc}
\usepackage[T1]{fontenc}
\usepackage{hyperref}
\title{Malinaltepec Me'phaa Language Resources for K-12 Curriculum Development}
\author{Compiled by Claude Sonnet 5 (claude-sonnet-5) for Andr\'as Kornai}
\date{}
\begin{document}
\maketitle
\section{Introduction}
Malinaltepec Me'phaa (Tlapanec; ISO 639-3 tcf) is an Otomanguean language of Guerrero, Mexico, with roughly 37,500 speakers according to search results. Related Tlapanec varieties have separate codes (e.g.\ tpx); do not mix them up.
\section{Primary Educational Resources}
No primers, textbooks or school programmes were identified in search results. A writing guide, \emph{C\'omo se escribe el tlapaneco} (Malinaltepec, Editorial Cuajimalpa, 1988), was named in a Wikipedia citation, but I saw no URL for it.
\section{Linguistic and Reference Materials}
Wikipedia overview of Tlapanec (cites Su\'arez 1988 \emph{Tlapaneco de Malinaltepec}, El Colegio de M\'exico; Carrasco Z\'u\~niga 1983, UNAM; Weathers and Weathers 1984 phonology sketch): \url{https://www.wikipedia.com/wiki/ISO_639:tpx}

CIESAS master's thesis on alignment in the Me'phaa of Malinaltepec (Abad Navarro Solano, 2012): \url{https://ciesas.repositorioinstitucional.mx/jspui/bitstream/1015/805/1/TE%20N.S.%202012%20Abad%20Navarro%20Solano.pdf}

Endangered Languages Alliance page on Tlapanec: \url{https://www.elalliance.org/languages/tlapanec}

No specific Malinaltepec dictionary was identified.
\section{Religious and Community Resources}
Background article on the Tlapaneco (Me'phaa) people: \url{https://www.indigenousmexico.org/articles/the-defiant-tlapaneca-mphaa-of-guerrero}
\section{Digital and Audio Resources}
None found in search results.
\section{Curriculum Development Context}
The language is not known to be taught in schools in the sources found. Linguistic descriptions exist, but pedagogical materials need further search (for example SIL Mexico and INALI catalogues, which I did not confirm).
\end{document}